\section*{Hoofdstuk \ref{ch:g12:age-of-the-earth} --- Radiometrische datering en de ouderdom van de Aarde}

\begin{solution}{exo:g12:age-of-the-earth:1}
(a) een \omterm{def:g12:age-of-the-earth:closed-system}{gesloten systeem}; (b) een \omterm{def:g12:age-of-the-earth:isochron}{isochroon}; (c) de \omterm{def:g12:age-of-the-earth:age}{ouderdom van de Aarde}.
\end{solution}

\begin{solution}{exo:g12:age-of-the-earth:2}
Plagioklaas (0.35), hoornblende (1.20), heel gesteente (3.10), K-veldspaat (6.40),
biotiet (24.0). Het biotiet ligt het hoogst: het bevat het meeste rubidium-87 voor
zijn strontium-86, dus heeft zijn rubidium het meeste strontium-87 gemaakt sinds het
graniet ontstond.
\end{solution}

\begin{solution}{exo:g12:age-of-the-earth:3}
$T = \ln 2/\lambda = 0.6931/\num{1.3972e-11} \approx \num{4.96e10}$ jaar,
49.6 miljard jaar: $49.6/4.54 \approx 11$ keer de \omterm{def:g12:age-of-the-earth:age}{ouderdom van de Aarde}.
\end{solution}

\begin{solution}{exo:g12:age-of-the-earth:4}
$\lambda t = \num{1.3972e-11} \times \num{3.10e8} \approx 0.004331$, en
$e^{0.004331} - 1 \approx 0.004341$. Voor 480 miljoen jaar is $\lambda t
\approx 0.006707$ en $e^{\lambda t} - 1 \approx 0.006729$. Beide liggen dicht bij
$\lambda t$ zelf, omdat $\lambda t$ klein is.
\end{solution}

\begin{solution}{exo:g12:age-of-the-earth:5}
Phillips (1860), Kelvin (1862), Joly (1899), Holmes (1927), Patterson
(1956). $4540/100 \approx 45$ keer.
\end{solution}

\begin{solution}{exo:g12:age-of-the-earth:6}
$0.7080 + 10.0 \times 0.004341 \approx 0.7514$ vandaag; 0.7080 toen het graniet
ontstond, zoals elk mineraal van het graniet.
\end{solution}

\begin{solution}{exo:g12:age-of-the-earth:7}
(1) $m = (0.77525 - 0.71806)/(9.40 - 0.90) = 0.05719/8.50 \approx 0.006728$;
$t = \ln(1.006728)/\num{1.3972e-11} \approx \num{4.80e8}$ jaar, 480 miljoen
jaar. (2) $0.71806 + 1.70 \times 0.006728 \approx 0.72950$ en $0.71806 + 4.20
\times 0.006728 \approx 0.74632$: beide op de rechte. (3) Ja: het graniet doorsnijdt
de schist, dus het is jonger, en 310 is minder dan 480 miljoen jaar.
\end{solution}

\begin{solution}{exo:g12:age-of-the-earth:8}
Al zijn punten zouden op dezelfde plaats vallen: één punt bepaalt geen rechte en
geen helling. Het biotiet ligt ver rechts van de andere; een kleine fout op
een punt verandert de helling van de rechte door dat punt en het plagioklaas dan
veel minder dan bij twee punten dicht bij elkaar.
\end{solution}

\begin{solution}{exo:g12:age-of-the-earth:9}
(1) Zijn rubidium is onveranderd, dus het houdt zijn $x$; het heeft minder
strontium-87, dus zijn $y$ daalt: het punt zakt recht naar beneden. (2) De rechte
door dat punt en het plagioklaas is minder steil, en geeft een te jonge ouderdom.
(3) Het biotiet ligt niet meer op de rechte door de andere mineralen: de
punten liggen niet op één lijn, een teken dat één monster geen \omterm{def:g12:age-of-the-earth:closed-system}{gesloten systeem} is geweest.
\end{solution}

\begin{solution}{exo:g12:age-of-the-earth:10}
(1) $4540/96 \approx 47$ keer. (2) Als er zout uit de zee verdwijnt, bijvoorbeeld
begraven in sedimenten, bevat de zee minder zout dan de rivieren in totaal aanvoerden:
het aanwezige zout delen door de jaarlijkse aanvoer geeft dan een tijd die te
kort is.
\end{solution}

\begin{solution}{exo:g12:age-of-the-earth:11}
$m = (34.86 - 10.34)/(50.28 - 9.46) = 24.52/40.82 \approx 0.601$, tussen
de hellingen voor 4.55 (0.590) en 4.60 miljard jaar (0.610): ongeveer 4.58
miljard jaar.
\end{solution}

\begin{solution}{exo:g12:age-of-the-earth:12}
(1) Het basalt, met de lagere verhouding, uit de mantel; het graniet, met de
hogere, minstens deels uit gesmolten korst. (2) Rubidium gaat naar de
smelt wanneer de mantel smelt, dus is het in de korst geconcentreerd, die
in de loop van de tijden meer strontium-87 heeft gemaakt. (3) Nee: mineralen die
uit één magma kristalliseren, delen zijn beginverhouding, en deze twee gesteenten niet.
\end{solution}

\begin{solution}{exo:g12:age-of-the-earth:13}
(1) Rubidium-87: $2^{-4.54/49.61} \approx 0.94$, ongeveer \qty{94}{\percent}
over; uranium-238: $2^{-4.54/4.47} \approx 0.49$, ongeveer de helft. (2) $m =
e^{\lambda t} - 1 \approx \num{1.3972e-11} \times 10^6 \approx
\num{1.4e-5}$: zelfs het biotiet van het graniet van Varden ($x = 24$) zou
maar $0.0003$ boven de beginverhouding uitstijgen. Er is zo weinig strontium-87 gevormd
dat de helling nauwelijks van nul te onderscheiden is: een trage klok past bij oude gesteenten.
\end{solution}

\begin{solution}{exo:g12:age-of-the-earth:14}
(1) $4.438 + 0.6026 \times 19.0 \approx 15.89$; $15.5$ ligt $0.39$ lager,
ongeveer \qty{2.4}{\percent}: het sediment ligt dicht bij de rechte. (2) 0.6026
ligt tussen 0.590 (4.55) en 0.610 (4.60): ongeveer 4.58 miljard jaar, binnen
het bereik van Patterson, van 4.48 tot 4.62 miljard jaar.
\end{solution}

\begin{solution}{exo:g12:age-of-the-earth:15}
(1) De voorraad verliest dan netto maar $46.7 - 24.4 = \qty{22.3}{TW}$, dus
houdt dezelfde voorraad $46.7/22.3 \approx 2.1$ keer langer stand. (2) $96 \times 2.1 \approx
200$ miljoen jaar, nog altijd meer dan twintig keer te kort. (3)
\omterm{def:g12:internal-heat:convection}{Mantelconvectie}, die warmte van diep binnenin naar boven voert en de gradiënt
dicht bij het oppervlak veel langer hoog houdt dan geleiding alleen.
\end{solution}

\begin{solution}{pb:g12:age-of-the-earth:1}
\textbf{1.} Het biotiet ($x = 24.0$).

\textbf{2.} Ze kristalliseerden uit één magma, en mineralen die uit één smelt
kristalliseren, nemen het strontium op met de isotopenverhouding van die smelt.

\textbf{3.} Strontium-86 is stabiel en wordt door geen enkel verval gemaakt: in een \omterm{def:g12:age-of-the-earth:closed-system}{gesloten systeem}
verandert zijn aantal niet, dus is het een vaste maatstaf.

\textbf{4.} Strontium-87, gemaakt door het verval van rubidium-87; rubidium-87 is
afgenomen en strontium-86 is gelijk gebleven.

\textbf{5.} Een horizontale lijn bij $y = 0.7080$: dezelfde verhouding, verschillende $x$.

\textbf{6.} $(0.81218 - 0.70952)/(24.0 - 0.35) = 0.10266/23.65 \approx
0.004341$.

\textbf{7.} $(0.72146 - 0.70952)/(3.10 - 0.35) = 0.01194/2.75 \approx
0.004342$.

\textbf{8.} Ze komen overeen tot op 1 op 4000. De eerste is betrouwbaarder:
haar twee punten liggen ver uit elkaar, dus de afronding van de verhoudingen telt minder.

\textbf{9.} $0.70952 - 0.35 \times 0.004341 \approx 0.7080$.

\textbf{10.} $0.7080 + 6.40 \times 0.004341 \approx 0.73578$: precies de
gemeten waarde.

\textbf{11.} $e^{\lambda t} = 1 + m = 1.004341$.

\textbf{12.} $\lambda t = \ln 1.004341 \approx 0.004331$.

\textbf{13.} $t = 0.004331/\num{1.3972e-11} \approx \num{3.10e8}$ jaar.

\textbf{14.} Ja: het graniet doorsnijdt de schist, dus volgens het principe van
doorsnijding is het jonger, en 310 is minder dan 480 miljoen
jaar.

\textbf{15.} 310 miljoen jaar ligt tussen 358.9 en 298.9 miljoen jaar:
het Carboon.

\textbf{16.} De verhouding $^{87}\mathrm{Sr}/^{86}\mathrm{Sr}$ van het magma toen
het graniet kristalliseerde.

\textbf{17.} Het graniet: $0.7080 > 0.7035$.

\textbf{18.} Rubidium gaat naar de smelt wanneer de mantel smelt, dus is het
in de korst geconcentreerd; met meer rubidium heeft de korst meer
strontium-87 gemaakt.

\textbf{19.} Dat het magma minstens voor een deel uit gesmolten \omterm{def:g12:continental-crust:continental-crust}{continentale
korst} kwam, en niet alleen uit de mantel.

\textbf{20.} Het graniet van Varden is \textbf{310 miljoen jaar} oud.
\end{solution}
